% Cote 119: offprints, reports, courses, CV and letters (typed and printed part): archive cote 119, site page 203 (PDF page 204).
% typescript. Draft compared: none.
% AI-transcribed by Claude Opus 5.5 (claude-opus-5-5, adaptive thinking, effort high), Message Batches API, batch r119_full (msgbatch_01DwZNm5Y7YsKBWyAN4sFs8w), 2026-11-13; not yet checked by eye.
\scanpage{119/203}
\mspage{- 8 -}

polynômes caractéristiques de l'endomorphisme de Frobenius opérant sur la
``cohomologie à support propre'' de la variété envisagée. Bien au delà d'une
simple question de rationalité, ces résultats ouvrent la voie à une \uline{approche}
cohomologique systématique d'invariants arithmétiques subtils comme les fonctions
\hand{$\zeta$} et L des variétés, et l'interprétation en termes arithmétiques de théo-%
rèmes tels que les théorèmes de dualité (démontrés à l'heure actuelle) et de
Lefschetz pour les sections hyperplanes (non démontrés encore en car. $>$ 0). Il
y a là un champ d'étude immense, qui par la nature des choses devrait se trouver,
tôt ou tard, centré sur la notion de ``motif'' (dénominateur commun des divers
types de cohomologie qu'on sait attacher à une variété algébrique) - mais qui
probablement ne sera jamais exploré jusqu'au bout, l'heure de ce genre d'inves-%
tigations étant déjà passée (même si rares sont ceux qui en ont pris conscience).

\uline{Références} : J.P. Serre, A. Weil, B. Dwork, J. Tate, M. Artin, P. Deligne ...

10. \uline{Géométrie Algébrique} ([12] à [19], [15 bis] à [20 bis]).

C'est dans cette direction que mon influence a été la plus directe et
la plus profonde, puisque c'est dans cette optique que se placent pour l'essen-%
tiel mes travaux depuis 1959. Voici les thèmes principaux sous lesquels on peut
placer mes contributions:

a) \uline{Travail de fondement}: Il s'agissait de dégager un cadre suffisamment vaste
pour servir de fondement commun à la géométrie algébrique habituelle (y com-%
pris celle développée par des auteurs comme A. Weil, O. Zariski, C. Chevalley,
J.P. Serre sur des corps de base quelconques) et à l'arithmétique. C'est fait
pour l'essentiel dans [15, Chap. I,II et des parties des Ch. III et IV], avec
l'introduction et l'étude de la \uline{notion de schéma}. Des généralisations ont été
développées par la suite, dans le même esprit, avec les schémas formels [15,
Chap. I, par. 10], la théorie des ``algebraic spaces'' de M. Artin (cf. Knutson
[*]), les ``algebraic stacks'' ou ``multiplicités algébriques'' de P. Deligne et
D. Mumford (*), des ``schémas relatifs'' de la thèse de M. Hakim [*] (en attendant
les ``multiplicités formelles'' et les ``multiplicités algébriques relatives'' sur
des topos annelés généraux,etc). Ces généralisations montrent la part concep-%
tuelle importante qui revient, dans le langage des schémas, à la notion générale
de la localisation, c'est à dire à celle de \uline{topos} (dont il a été question au
par. 3). Les fondements développés dans [15] et [16] sont aujourd'hui le ``pain
quotidien'' de la grande majorité des géomètres algébristes, et leur importance
a été soulignée à de nombreuses occasions par des mathématiciens aussi divers

\margin{203}
